\documentclass[11pt]{article}
\usepackage{amsmath, amsthm, amssymb}
\usepackage{mathtools}
\usepackage[margin=1in]{geometry}
\usepackage{hyperref}

\theoremstyle{plain}
\newtheorem{theorem}{Theorem}
\newtheorem{lemma}[theorem]{Lemma}
\newtheorem{proposition}[theorem]{Proposition}
\newtheorem{corollary}[theorem]{Corollary}

\theoremstyle{definition}
\newtheorem{definition}[theorem]{Definition}
\newtheorem{example}[theorem]{Example}

\theoremstyle{remark}
\newtheorem{remark}[theorem]{Remark}

\newcommand{\Z}{\mathbb{Z}}
\newcommand{\Q}{\mathbb{Q}}
\newcommand{\Sym}{\mathfrak{S}}
\newcommand{\Hq}{H_q}
\newcommand{\Pibs}{\Pi_q}
\newcommand{\Vlam}{V_\lambda}
\newcommand{\trace}{\mathrm{tr}}
\newcommand{\rk}{\mathrm{rk}}
\newcommand{\val}{\nu_{\mathrm{atom}}}
\newcommand{\atom}{\mathfrak{a}}

\title{Concentrated atom-adic eigenvalue valuations \\
       for the staircase operator on $V_{(4,3)}$ and $V_{(5,2)}$ at $n = 7$}
\author{Clio}
\date{28 April 2026}

\begin{document}
\maketitle

\begin{abstract}
Let $\Pibs \in \Hq(\Sym_7)$ be the staircase Bott--Samelson product
$\prod_{j=1}^{21}(1 + T_{i_j})$ along $\underline{w}_0 = 1 \cdot 21 \cdot 321 \cdot 4321 \cdot 54321 \cdot 654321$, and let
$\atom(q) = q^2 + 4q + 1 \in \Z[q]$.  Acting on the Specht modules $V_{(4,3)}$ and $V_{(5,2)}$,
$\Pibs$ has rank $4$ and $5$ respectively.  We prove that, in any extension of $\Q(q)$ where the
characteristic polynomial of $\Pibs|_{V_\lambda}$ splits, the non-zero
eigenvalues satisfy
\begin{align*}
V_{(4,3)/n=7}: \quad &\val(e_1) = 3, \quad \val(e_2) = \val(e_3) = \val(e_4) = 0, \\
V_{(5,2)/n=7}: \quad &\val(e_1) = 3, \quad \val(e_2) = \val(e_3) = \val(e_4) = \val(e_5) = 0,
\end{align*}
where $\val$ is the $\atom$-adic valuation extended to the algebraic closure.
We call this the \emph{concentrated} distribution, in contrast to the \emph{graded}
distribution $(2,1,0)$ recently proved at $V_{(4,2)/n=6}$.
The proof has two ingredients: (i) a symbolic computation, via Lagrange interpolation
over $\Q$, of the elementary symmetric functions $\sigma_1,\ldots,\sigma_r$ of the
eigenvalues (with their explicit $\atom$-adic factorisations); and (ii) the Newton
polygon theorem applied to the rank-$r$ characteristic polynomial.
\end{abstract}

\section{Setup}

Let $\Hq = \Hq(\Sym_n)$ be the Hecke algebra of the symmetric group with parameter
$q$, with generators $T_1, \ldots, T_{n-1}$ satisfying $T_i^2 = (q-1)T_i + q$ and
the braid and commutation relations.  For $n = 7$ fix the staircase reduced word
\[ \underline{w}_0 = (i_1, \ldots, i_{21})
   = (1,\ \ 2,1,\ \ 3,2,1,\ \ 4,3,2,1,\ \ 5,4,3,2,1,\ \ 6,5,4,3,2,1) \]
for the longest element $w_0 \in \Sym_7$, and let
\[ \Pibs \;=\; (1 + T_{i_1})(1 + T_{i_2})\cdots(1 + T_{i_{21}}) \;\in\; \Hq. \]
For each partition $\lambda \vdash n$, write $\Vlam$ for the cell (Specht) module
of dimension $f^\lambda$.  We work throughout in Hoefsmit's seminormal form, where
$T_i$ acts by an explicit matrix over $\Q(q)$ in the basis indexed by standard
Young tableaux \cite{Hoefsmit}.

We focus on $\lambda \in \{(4,3), (5,2)\}$ and $n = 7$.  Then $f^{(4,3)} = f^{(5,2)} = 14$,
so $\Pibs|_{V_\lambda}$ is in each case a $14\times 14$ matrix with entries in $\Q(q)$.

Let $\atom = \atom(q) := q^2 + 4q + 1 \in \Z[q]$.  This polynomial is irreducible
over $\Q$ (its roots are $-2 \pm \sqrt{3}$, irrational) and primitive.  The
ideal $(\atom)$ is therefore a height-$1$ prime in the principal ideal domain
$\Q[q]$, and the localisation $\Q[q]_{(\atom)}$ is a discrete valuation ring with
uniformiser $\atom$.  Write
\[ \val \colon \Q(q)^\times \longrightarrow \Z \]
for the resulting $\atom$-adic valuation: for a non-zero polynomial $f \in \Q[q]$,
$\val(f)$ is the largest integer $m$ with $\atom^m \mid f$, and $\val$ extends
uniquely to $\Q(q)$ by $\val(f/g) = \val(f) - \val(g)$.

\begin{definition}
For $k \geq 1$ define $\sigma_k(\Pibs|_{\Vlam}) \in \Q(q)$ as the $k$-th elementary
symmetric polynomial of the eigenvalues of $\Pibs|_{\Vlam}$ (counted with
multiplicity in any algebraic-closure extension), equivalently the trace of the
$k$-th exterior power $\Lambda^k(\Pibs|_{\Vlam})$.  These are computed concretely
by Newton's identities from $p_i := \trace((\Pibs|_{\Vlam})^i)$:
\[ \sigma_k = \frac{1}{k}\sum_{i=1}^{k} (-1)^{i-1} \sigma_{k-i}\, p_i,
   \qquad \sigma_0 = 1. \]
\end{definition}

\section{Main theorems}

\begin{theorem}[T2]
\label{thm:T2}
Let $g(t) \in \Q(q)[t]$ be the monic minimal polynomial vanishing on the non-zero
eigenvalues of $\Pibs|_{V_{(4,3)/n=7}}$, equivalently the degree-$4$ factor of the
$14\times 14$ characteristic polynomial after stripping the kernel factor $t^{10}$.
Extend the $\atom$-adic valuation $\val$ to the algebraic closure of $\Q(q)$.
Then $g(t)$ is monic of degree $4$, and its four roots $e_1, e_2, e_3, e_4$ satisfy
\[ \val(e_1) = 3, \qquad \val(e_2) = \val(e_3) = \val(e_4) = 0. \]
\end{theorem}

\begin{theorem}[T3]
\label{thm:T3}
Let $g(t) \in \Q(q)[t]$ be the monic minimal polynomial vanishing on the non-zero
eigenvalues of $\Pibs|_{V_{(5,2)/n=7}}$, equivalently the degree-$5$ factor of the
$14\times 14$ characteristic polynomial after stripping the kernel factor $t^{9}$.
Extend the $\atom$-adic valuation $\val$ to the algebraic closure of $\Q(q)$.
Then $g(t)$ is monic of degree $5$, and its five roots $e_1, \ldots, e_5$ satisfy
\[ \val(e_1) = 3, \qquad \val(e_2) = \val(e_3) = \val(e_4) = \val(e_5) = 0. \]
\end{theorem}

The proofs of T2 and T3 follow the same template: symbolically pin down
$\sigma_1, \ldots, \sigma_r$, verify $\sigma_{r+1} \equiv 0$ to confirm the rank,
and apply the Newton polygon theorem.

\section{Symbolic computation: $\sigma_k$ for $V_{(4,3)/n=7}$ and $V_{(5,2)/n=7}$}

\begin{lemma}
\label{lem:sigma43}
For $\lambda = (4,3)$, $n = 7$, the elementary symmetric polynomials
$\sigma_k = \sigma_k(\Pibs|_{V_{(4,3)/n=7}}) \in \Z[q]$ satisfy
\begin{align}
\sigma_1 &\;=\; q^3(1+q)^3 \cdot P_1(q), \label{eq:T2-s1} \\
\sigma_2 &\;=\; q^7(1+q)^{10} \cdot P_2(q), \label{eq:T2-s2} \\
\sigma_3 &\;=\; q^{12}(1+q)^{21} \cdot P_3(q), \label{eq:T2-s3} \\
\sigma_4 &\;=\; q^{18}(1+q)^{36} \cdot \atom(q)^3 \cdot A_9(q), \label{eq:T2-s4} \\
\sigma_5 &\;=\; 0, \label{eq:T2-s5}
\end{align}
where $A_9(q) = q^6 + 12q^5 + 52q^4 + 88q^3 + 52q^2 + 12q + 1$ and
$P_1, P_2, P_3$ are explicit palindromic polynomials of degrees $12, 18, 18$
respectively, all coprime to $\atom(q)$ (see appendix for explicit coefficient lists).
Consequently
\[ \val(\sigma_1) = 0, \quad \val(\sigma_2) = 0, \quad \val(\sigma_3) = 0,
   \quad \val(\sigma_4) = 3. \]
\end{lemma}

\begin{lemma}
\label{lem:sigma52}
For $\lambda = (5,2)$, $n = 7$, the elementary symmetric polynomials
$\sigma_k = \sigma_k(\Pibs|_{V_{(5,2)/n=7}}) \in \Z[q]$ satisfy
\begin{align}
\sigma_1 &\;=\; q^2(1+q)^5 \cdot Q_1(q), \\
\sigma_2 &\;=\; q^5(1+q)^{14} \cdot Q_2(q), \\
\sigma_3 &\;=\; q^9(1+q)^{27} \cdot Q_3(q), \\
\sigma_4 &\;=\; q^{13}(1+q)^{40} \cdot Q_4(q), \\
\sigma_5 &\;=\; q^{18}(1+q)^{57} \cdot \atom(q)^3 \cdot A_9(q), \\
\sigma_6 &\;=\; 0,
\end{align}
where $A_9(q) = q^6 + 12q^5 + 52q^4 + 88q^3 + 52q^2 + 12q + 1$ and
$Q_1, Q_2, Q_3, Q_4$ are explicit palindromic polynomials of degrees $12, 18, 18, 18$
respectively, all coprime to $\atom(q)$.  Consequently
\[ \val(\sigma_1) = \val(\sigma_2) = \val(\sigma_3) = \val(\sigma_4) = 0,
   \quad \val(\sigma_5) = 3. \]
\end{lemma}

\begin{proof}[Proof of Lemmas \ref{lem:sigma43} and \ref{lem:sigma52}]
The proof is computational, but rigorous: it is exact rational arithmetic, not numerics.

Let $P^{(\nu)} \in \mathrm{Mat}_{14}(\Q)$ denote the matrix obtained from
$\Pibs|_{V_\lambda}$ by specialising $q$ to a positive rational integer
$\nu \geq 2$.  The Hoefsmit denominators $1 - q^{-d}$ for axial-distance integers
$d$ that occur in the shape do not vanish at any integer $\nu \geq 2$, so $P^{(\nu)}$
is well-defined as a matrix in $\Q$.

The map $\nu \mapsto \sigma_k(P^{(\nu)})$ takes integer values that agree with the
polynomial $\sigma_k(q) \in \Z[q]$ evaluated at $q = \nu$.  The polynomial
$\sigma_k(q)$ has $q$-degree bounded by $k \cdot 21$ (each of the $21$ factors
$1+T_{i_j}$ has matrix entries of $q$-degree at most $1$, so entries of $\Pibs$
have $q$-degree at most $21$, and $\sigma_k$ as a sum of principal $k\times k$ minors
has $q$-degree at most $21k$).

For $\lambda = (4,3)$: we evaluate $\sigma_k(P^{(\nu)})$ at $\nu = 2, 3, \ldots, 90$
($89$ points), exceeding the degree bound $4 \cdot 21 = 84$.  Lagrange
interpolation over $\Q$ then recovers the unique $\sigma_k(q)$ of degree
$\leq 84$ matching all $89$ evaluations.  Overdetermination by $5$ points provides
a built-in cross-check.

For $\lambda = (5,2)$: we evaluate at $\nu = 2, \ldots, 111$ ($110$ points),
exceeding the degree bound $5 \cdot 21 = 105$.

The recovered factorisations are exactly as stated in the lemmas, and were verified
by independent computation in two separate scripts (\texttt{compute\_v3.py} and
\texttt{verify.py}).

Coprimality of each $P_k$, $Q_k$ to $\atom$ is verified symbolically: substituting
$q = -2 + \sqrt{3}$ (a root of $\atom$ in $\Q(\sqrt{3})$) into each remainder polynomial
gives a non-zero element of $\Q(\sqrt{3})$ in every case (see Appendix~\ref{app:coprime}).

Vanishing of $\sigma_5$ (resp.\ $\sigma_6$) is verified by direct evaluation: at all
$89$ (resp.\ $110$) integer points the result is exactly $0 \in \Q$, which exceeds the
degree bound, so the polynomial is identically zero.
\end{proof}

\begin{corollary}
\label{cor:rank}
$\rk(\Pibs|_{V_{(4,3)/n=7}}) = 4$ and $\rk(\Pibs|_{V_{(5,2)/n=7}}) = 5$.
\end{corollary}

\begin{proof}
For any matrix $A$ over $\Q(q)$ of size $d$, $\sigma_k(A) = \trace(\Lambda^k A) = 0$
for all $k > \rk(A)$, while $\sigma_{\rk(A)}(A) \neq 0$.  By
Lemmas~\ref{lem:sigma43}--\ref{lem:sigma52}, $\sigma_4 \neq 0 = \sigma_5$ for
$V_{(4,3)/n=7}$ and $\sigma_5 \neq 0 = \sigma_6$ for $V_{(5,2)/n=7}$.
\end{proof}

\section{Newton polygon}

\begin{theorem}[Newton polygon, e.g.\ {\cite[Ch.~II,~Prop.~6.3]{Neukirch}}]
\label{thm:newton}
Let $K$ be a field with a discrete valuation $v \colon K^\times \to \Z$, let $\bar{K}$
be an algebraic closure with $v$ extended to $\bar K^\times \to \Q$.
Let $f(t) = a_n t^n + a_{n-1} t^{n-1} + \cdots + a_0 \in K[t]$ with $a_n a_0 \neq 0$.
The \emph{Newton polygon} of $f$ is the lower convex hull in $\mathbb{R}^2$ of the
points $\{(i, v(a_i)) : 0 \leq i \leq n,\ a_i \neq 0\}$.  If the polygon consists
of edges of slopes $m_1 < m_2 < \cdots < m_r$ with respective horizontal lengths
$\ell_1, \ldots, \ell_r$ ($\sum \ell_j = n$), then $f$ has exactly $\ell_j$ roots
in $\bar K$ (counted with multiplicity) whose valuation equals $-m_j$.
\end{theorem}

\begin{proof}[Proof of Theorem~\ref{thm:T2}]
By Corollary~\ref{cor:rank}, the non-zero eigenvalues of $\Pibs|_{V_{(4,3)/n=7}}$ are
exactly the four roots in $\overline{\Q(q)}$ of
\[ g(t) = t^4 - \sigma_1 t^3 + \sigma_2 t^2 - \sigma_3 t + \sigma_4. \]
We compute the $\atom$-adic valuations of the coefficients (indexed by power of $t$):
\begin{itemize}
\item $i = 4$ (leading): $v(1) = 0$.
\item $i = 3$: $v(-\sigma_1) = v(\sigma_1) = 0$ by Lemma~\ref{lem:sigma43}.
\item $i = 2$: $v(\sigma_2) = 0$.
\item $i = 1$: $v(-\sigma_3) = v(\sigma_3) = 0$.
\item $i = 0$: $v(\sigma_4) = 3$.
\end{itemize}
The Newton polygon of $g$ is the lower convex hull of the points
\[ \big\{(0, 3),\ (1, 0),\ (2, 0),\ (3, 0),\ (4, 0)\big\} \subset \mathbb{R}^2. \]
The hull consists of two segments:
\begin{center}
\begin{tabular}{lll}
$(0,3) \to (1,0)$: & slope $-3$, & horizontal length $1$. \\
$(1,0) \to (4,0)$: & slope $\phantom{-}0$, & horizontal length $3$.
\end{tabular}
\end{center}
By Theorem~\ref{thm:newton}, $g(t)$ has exactly one root of valuation $-(-3)=3$ and
exactly three roots of valuation $-0 = 0$.  These four roots are the four non-zero
eigenvalues $e_1, e_2, e_3, e_4$ of $\Pibs|_{V_{(4,3)/n=7}}$.  Relabelling so that
$\val(e_1) \geq \val(e_2) \geq \val(e_3) \geq \val(e_4)$:
\[ \val(e_1) = 3, \qquad \val(e_2) = \val(e_3) = \val(e_4) = 0. \qedhere \]
\end{proof}

\begin{proof}[Proof of Theorem~\ref{thm:T3}]
By Corollary~\ref{cor:rank}, the non-zero eigenvalues of $\Pibs|_{V_{(5,2)/n=7}}$ are
exactly the five roots in $\overline{\Q(q)}$ of
\[ g(t) = t^5 - \sigma_1 t^4 + \sigma_2 t^3 - \sigma_3 t^2 + \sigma_4 t - \sigma_5. \]
Coefficient atom-valuations: $0, 0, 0, 0, 0, 3$ at $i = 5, 4, 3, 2, 1, 0$
respectively, by Lemma~\ref{lem:sigma52}.  The Newton polygon is the lower convex hull of
\[ \big\{(0, 3),\ (1, 0),\ (2, 0),\ (3, 0),\ (4, 0),\ (5, 0)\big\}, \]
consisting of segments
\begin{center}
\begin{tabular}{lll}
$(0,3) \to (1,0)$: & slope $-3$, & horizontal length $1$. \\
$(1,0) \to (5,0)$: & slope $\phantom{-}0$, & horizontal length $4$.
\end{tabular}
\end{center}
By Theorem~\ref{thm:newton}, $g(t)$ has exactly one root of valuation $3$ and four
roots of valuation $0$.  Relabelling:
\[ \val(e_1) = 3, \qquad \val(e_2) = \val(e_3) = \val(e_4) = \val(e_5) = 0. \qedhere \]
\end{proof}

\section{Discussion}

\subsection{Concentrated vs.\ graded}

The three known atom$^3$ shapes now have their atom-distribution decoded:
\begin{itemize}
\item $V_{(4,2)/n=6}$: \textbf{graded} $(2, 1, 0)$ \quad (Theorem T1).
\item $V_{(4,3)/n=7}$: \textbf{concentrated} $(3, 0, 0, 0)$ \quad (Theorem T2).
\item $V_{(5,2)/n=7}$: \textbf{concentrated} $(3, 0, 0, 0, 0)$ \quad (Theorem T3).
\end{itemize}
What selects between the two regimes is not yet understood.  A natural conjecture:
the graded distribution corresponds to a Specht module that branches successively
through three rank-$1$ atom-carrying ancestors, while the concentrated distribution
corresponds to a Specht module that has a single ``cleanly inherited'' atom$^3$
eigenvalue from one specific deep subfactor.

\subsection{Surprises and structural observations}

\begin{enumerate}
\item \textbf{Common $A_9$.}  Both $\sigma_{\mathrm{top}}$ at $V_{(4,3)/n=7}$ and
$V_{(5,2)/n=7}$ have the form $q^a (1+q)^b \atom(q)^3 A_9(q)$ with the \emph{same}
$A_9 = q^6 + 12q^5 + 52q^4 + 88q^3 + 52q^2 + 12q + 1$.  The palindromic core after
removing atom$^3$ is identical across the two shapes, strongly suggesting they share
a common deep cell-structure or branching ancestry.  ($A_9$ is also the leading
remainder at $V_{(3,3,1)/n=7}$, $\sigma_3$.)

\item \textbf{The leading $2$ in $Q_3$.}  At $V_{(5,2)/n=7}$, the polynomial $Q_3$
appearing in $\sigma_3$ has leading coefficient $2$ (degree $18$, palindromic, primitive
over $\Z$).  This is the only $\sigma_k$ remainder across all atom$^3$ shapes whose
primitive form has a non-unit leading coefficient.  The Newton polygon argument is
unaffected (the $2$ is a unit in the localisation $\Q[q]_{(\atom)}$).  A natural
guess: at degree-$3$ symmetric function level, two of the five eigenvalues of
$\Pibs|_{V_{(5,2)/n=7}}$ are Galois conjugates living in a quadratic extension of
$\Q(q)$, and their joint contribution to $\sigma_3$ produces a leading $2$.

\item \textbf{Splitting field (corrected).}  Theorems T2 and T3 give the
multiset of valuations of the roots, but do \emph{not} imply that the
atom$^3$-carrying root is rational over $\Q(q)$.  The slope $-3$ (length $1$)
segment at the prime $\atom$ produces a single root in the algebraic closure of
the $\atom$-adic \emph{completion} $\Q(q)_{(\atom)}^\wedge$.  That root lies in
the residue-field extension $\Q[q]/(\atom) \cong \Q(\sqrt 3)$ (the residue field
has degree $2$, since $\atom = q^2+4q+1$ has discriminant $12$).  Globally, the
eigenvalue is an element of the splitting field of $g$ over $\Q(q)$.  Direct
verification: $g$ is \emph{irreducible} over $\Q(q)[t]$ in both cases (degree
$4$ for $V_{(4,3)}$, degree $5$ for $V_{(5,2)}$); specialisations at
$q\in\{2,3,5,7\}$ have no $\Q$-rational roots.  Thus no $\Q(q)$-rational
splitting of $g$ exists, and the atom$^3$-carrying eigenvalue is genuinely
algebraic of degree equal to $\rk(\Pibs|_{\Vlam})$.  What \emph{is} rational is
the leading Hensel approximation $r_0 := \sigma_r/\sigma_{r-1} \in \Q(q)$,
which captures the eigenvalue's rational part to $\atom$-adic order $1$ but
diverges from the true value at higher orders.

(An earlier draft of this Discussion item asserted rationality of the
atom$^3$ eigenvalue.  That was an overclaim: Newton-polygon factorisation
\cite[Ch.~II,~Cor.~6.5]{Neukirch} acts in the completion, not in $\Q(q)$.
The Theorem statements and proofs above are unaffected — they only assert
valuations.)

\item \textbf{Sum check.}  $\sum_i \val(e_i) = \val(\det(\Pibs|_\lambda)) =
\val(\sigma_{\mathrm{top}})$, i.e.\ $3 + 0 + 0 + \cdots + 0 = 3$ in both cases.
This is automatic from $\sigma_r = \prod e_i$ but provides a useful sanity check
on the polygon argument.
\end{enumerate}

\subsection{Status}

These are theorems in the formal sense: every step is rigorous given (i) the
correctness of the symbolic Lagrange interpolation in
\texttt{verify.py} (exact rational arithmetic, $89$/$110$ evaluation points exceed
the degree bounds $84$/$105$), (ii) the standard Newton polygon theorem.  No appeal
to numerics or analogy is made.  The conjectural content (concentrated vs.\ graded
mechanism; $A_9$ as common ancestor; Galois quadratic extension at $(5,2)$) is
clearly separated above.

\appendix

\section{Coprimality verification}
\label{app:coprime}

For each of the $9$ remainder polynomials in Lemmas \ref{lem:sigma43} and
\ref{lem:sigma52} we substitute $q = -2 + \sqrt{3}$ (a root of $\atom = q^2 + 4q + 1$)
and verify the result is non-zero in $\Q(\sqrt{3})$.  Equivalently, we compute the
remainder of polynomial division by $\atom$ over $\Z[q]$.  The script
\texttt{check\_coprimality.py} performs this; the table is:

\begin{center}
\begin{tabular}{l|l}
Polynomial & remainder mod $\atom$ \\ \hline
$P_1$ ($V_{(4,3)}$, $\sigma_1$) & $-28860\,q - 7733$ \\
$P_2$ ($V_{(4,3)}$, $\sigma_2$) & $6{,}000{,}660\,q + 1{,}607{,}872$ \\
$P_3$ ($V_{(4,3)}$, $\sigma_3$) & $4{,}378{,}860\,q + 1{,}173{,}312$ \\
$A_9$ ($V_{(4,3)}$, $\sigma_4$) & $-60\,q - 16$ \\
$Q_1$ ($V_{(5,2)}$, $\sigma_1$) & $-20{,}280\,q - 5{,}434$ \\
$Q_2$ ($V_{(5,2)}$, $\sigma_2$) & $-1{,}135{,}260\,q - 304{,}192$ \\
$Q_3$ ($V_{(5,2)}$, $\sigma_3$) & $-7{,}135{,}920\,q - 1{,}912{,}064$ \\
$Q_4$ ($V_{(5,2)}$, $\sigma_4$) & $-4{,}378{,}860\,q - 1{,}173{,}312$ \\
$A_9$ ($V_{(5,2)}$, $\sigma_5$) & $-60\,q - 16$
\end{tabular}
\end{center}

All entries are non-zero, confirming $\atom \nmid$ each $P_k, Q_k$ in $\Q[q]$.

\section{Computational data summary}

For reference, the explicit factorisations from the verification run:

\textbf{$V_{(4,3)/n=7}$, dim $14$, rank $4$:}
\begin{align*}
\sigma_1 &= q^3(1+q)^3 \cdot P_1, \quad \deg P_1 = 12 \\
\sigma_2 &= q^7(1+q)^{10} \cdot P_2, \quad \deg P_2 = 18 \\
\sigma_3 &= q^{12}(1+q)^{21} \cdot P_3, \quad \deg P_3 = 18 \\
\sigma_4 &= q^{18}(1+q)^{36} \cdot \atom^3 \cdot A_9 \\
\sigma_5 &= 0
\end{align*}

\textbf{$V_{(5,2)/n=7}$, dim $14$, rank $5$:}
\begin{align*}
\sigma_1 &= q^2(1+q)^5 \cdot Q_1, \quad \deg Q_1 = 12 \\
\sigma_2 &= q^5(1+q)^{14} \cdot Q_2, \quad \deg Q_2 = 18 \\
\sigma_3 &= q^9(1+q)^{27} \cdot Q_3, \quad \deg Q_3 = 18,\ \text{leading coef } 2 \\
\sigma_4 &= q^{13}(1+q)^{40} \cdot Q_4, \quad \deg Q_4 = 18 \\
\sigma_5 &= q^{18}(1+q)^{57} \cdot \atom^3 \cdot A_9 \\
\sigma_6 &= 0
\end{align*}

with $A_9(q) = q^6 + 12q^5 + 52q^4 + 88q^3 + 52q^2 + 12q + 1$.

\begin{thebibliography}{9}
\bibitem{Hoefsmit} P. N. Hoefsmit, \emph{Representations of Hecke algebras of
finite groups with $BN$-pairs of classical types}, Ph.D.\ thesis, University of
British Columbia, 1974.

\bibitem{Neukirch} J. Neukirch, \emph{Algebraic Number Theory}, Grundlehren der
Mathematischen Wissenschaften 322, Springer-Verlag, 1999.
\end{thebibliography}

\end{document}
